\expandafter\ifx\csname CAPETINCLUDEMODE\endcsname\relax
\documentclass[11pt,a4paper]{article}
% preamble.tex — packages et commandes communes a toutes les fiches CAPET SII IE
\input{glyphtounicode}
\pdfgentounicode=1
\usepackage{cmap}
\usepackage[T1]{fontenc}
\usepackage[utf8]{inputenc}
\usepackage[french]{babel}
\usepackage[a4paper, top=1.6cm, bottom=1.6cm, left=1.4cm, right=1.4cm]{geometry}
\setlength{\headheight}{14.5pt}
\usepackage{amsmath, amssymb, amsfonts}
\usepackage{siunitx}
\sisetup{locale=FR, output-decimal-marker={,}, inter-unit-product=\ensuremath{{}\cdot{}}}
\usepackage{array, booktabs, multirow, tabularx}
\usepackage{xcolor}
\usepackage[most, breakable]{tcolorbox}
\usepackage{multicol}
\usepackage{enumitem}
\usepackage{mathtools}
\usepackage{lmodern}
\usepackage{microtype}
\usepackage{fancyhdr}
\usepackage{tikz}
\usetikzlibrary{positioning,arrows.meta}
\usepackage{lastpage}

% ── Couleurs par domaine ──────────────────────────────────────────────
\definecolor{colElec}{RGB}{26, 82, 148}       % Électrotechnique — bleu foncé
\definecolor{colElecLight}{RGB}{210, 225, 245} % fond bleu clair
\definecolor{colAuto}{RGB}{34, 120, 66}        % Automatique — vert
\definecolor{colMeca}{RGB}{185, 90, 20}        % Mécanique — orange
\definecolor{colTherm}{RGB}{160, 40, 40}       % Thermique — rouge
\definecolor{colQT}{RGB}{232, 244, 233}        % fond question type (vert très léger)
\definecolor{colJury}{RGB}{255, 243, 205}      % fond jury (jaune ambre)
\definecolor{colJuryFrame}{RGB}{200, 130, 0}   % bordure jury
\definecolor{colGrayLight}{RGB}{245, 245, 247} % fond gris neutre
\definecolor{colSeparator}{RGB}{180, 180, 180} % gris séparateur

% ── Style tcolorbox de section principale ────────────────────────────
\tcbset{
  ficheSection/.style={
    breakable,
    enhanced,
    colback=#1!7!white,
    colframe=#1,
    colbacktitle=#1,
    coltitle=white,
    fonttitle=\bfseries\normalsize,
    left=4pt, right=4pt, top=3pt, bottom=4pt,
    before skip=6pt, after skip=6pt,
  }
}

% ── Style encadré question type ──────────────────────────────────────
\tcbset{
  questiontype/.style={
    breakable,
    enhanced,
    colback=colQT,
    colframe=colElec!55!black,
    colbacktitle=colElec!55!black,
    coltitle=white,
    fonttitle=\bfseries\small,
    left=4pt, right=4pt, top=3pt, bottom=3pt,
    before skip=5pt, after skip=5pt,
  }
}

% ── Style encadré jury ───────────────────────────────────────────────
\tcbset{
  jury/.style={
    breakable,
    enhanced,
    colback=colJury,
    colframe=colJuryFrame,
    fonttitle=\bfseries\normalsize,
    left=4pt, right=4pt, top=3pt, bottom=3pt,
  }
}

% ── Mise en page ─────────────────────────────────────────────────────
\pagestyle{fancy}
\fancyhf{}
\renewcommand{\headrulewidth}{0.4pt}
\renewcommand{\footrulewidth}{0.3pt}

\setlist[itemize]{leftmargin=1.2em, itemsep=1pt, topsep=2pt, parsep=0pt}
\setlist[enumerate]{leftmargin=1.2em, itemsep=3pt, topsep=2pt, parsep=0pt}

\setlength{\parskip}{2pt}
\setlength{\columnsep}{10pt}

\def\CAPETBEGINFILE{\begin{document}}
\def\CAPETENDFILE{\end{document}}
\else
\def\CAPETBEGINFILE{}
\def\CAPETENDFILE{}
\fi

\fancyhead[L]{\textbf{\textcolor{colElec}{T12}} - Machine synchrone et brushless (PMSM / BLDC)}
\fancyhead[R]{\small\textit{CAPET SII IE - Session 2026}}
\fancyfoot[C]{\small\thepage\,/\,\pageref{LastPage}}

\CAPETBEGINFILE

% -----------------------------------------------------------------------
%  BLOC EN-TeTE
% -----------------------------------------------------------------------

\begin{tcolorbox}[
  enhanced, breakable,
  colback=colElec,
  colframe=colElec,
  left=6pt, right=6pt, top=5pt, bottom=5pt,
  before skip=0pt, after skip=4pt
]
  {\color{white}\bfseries\large T12 \quad Machine synchrone et brushless (PMSM\,/\,BLDC)}\\[2pt]
  {\color{colElec!30!white}\small
    \textbf{Domaine :} electrotechnique - Machines electriques \quad
    \textbf{Presence :} \textbf{6/13} sujets (2013-2026) - theme \emph{regulier}
  }\\[1pt]
  {\color{colElec!20!white}\footnotesize
    Sous-themes : frequence d'alimentation et paires de peles - couplage etoile/triangle - diagramme de Kapp - servomoteur brushless (PMSM/BLDC) - autopilotage - rendement chaene de puissance
  }
\end{tcolorbox}

% -----------------------------------------------------------------------
%  1. RAPPELS DE COURS
% -----------------------------------------------------------------------

\begin{tcolorbox}[
  ficheSection=colElec,
  title={1. Rappels de cours},
]

% -- 1.1 Vitesse de synchronisme ---------------------------------------
\textbf{1.1 Vitesse de synchronisme}

\smallskip
\begin{center}
\renewcommand{\arraystretch}{1.5}
\begin{tabular}{lll}
\toprule
\textbf{Grandeur} & \textbf{Expression} & \textbf{Unite} \\
\midrule
Vitesse angulaire mecanique & $\omega_m = \dfrac{2\pi f}{p}$ & rades$^{-1}$ \\[6pt]
Vitesse de synchronisme & $N_s = \dfrac{60\,f}{p}$ & tremin$^{-1}$ \\[6pt]
Frequence imposee par la vitesse & $f = p \cdot \dfrac{N}{60}$ & Hz \\[4pt]
\bottomrule
\end{tabular}
\end{center}
$p$ = nombre de \textbf{paires} de peles ( nombre de peles).
Une machine - 28 peles a $p = 14$.

\medskip\hrule\vspace{4pt}

% -- 1.2 equation electrique et diagramme de Kapp ---------------------
\textbf{1.2 equation electrique - circuit equivalent monophase}

\smallskip
\[
  \vec{V}_s = \vec{E} + R_s \vec{I}_s + jX_s \vec{I}_s
  \qquad;\qquad
  X_s = L_s \omega
\]
En convention moteur avec $R_s \approx 0$ :
\[
  E = \sqrt{V_s^2 + (X_s I_s)^2} \quad \text{(si } \cos\varphi = 1 \text{)}
  \qquad;\qquad
  P_{max} = \frac{3\,V_s E}{X_s} \quad (\delta = 90^\circ)
\]
L'angle de charge $\delta$ : en moteur $\delta > 0$ (rotor en avance) ; en generateur $\delta < 0$.

\medskip\hrule\vspace{4pt}

% -- 1.3 Servomoteur PMSM/BLDC ----------------------------------------
\textbf{1.3 Servomoteur brushless (PMSM / BLDC) - comparaison MCC}

\smallskip
\begin{center}
\renewcommand{\arraystretch}{1.25}
\begin{tabular}{llll}
\toprule
\textbf{Critere} & \textbf{MCC} & \textbf{PMSM / BLDC} & \textbf{MAS} \\
\midrule
Entretien & Balais - changer & Aucun (pas de balais) & Faible \\
Commutation & Mecanique (collecteur) & electronique (onduleur) & - \\
Densite de puissance & Moyenne & elevee & Moyenne \\
Coet & Moyen & eleve & Faible \\
Capteur position & Non requis & Resolveur / Encodeur / Hall & Non requis \\
\bottomrule
\end{tabular}
\end{center}

\medskip

\textbf{Servomoteur vs moteur :} un servomoteur est un moteur associe - un capteur de retour (resolveur, encodeur, capteurs Hall) et - une boucle fermee d'asservissement en position, vitesse ou couple.

\medskip\hrule\vspace{4pt}

% -- 1.4 Couplage etoile/triangle -------------------------------------
\textbf{1.4 Couplage etoile / triangle}

\smallskip
\begin{center}
\renewcommand{\arraystretch}{1.25}
\begin{tabular}{lll}
\toprule
\textbf{Couplage} & \textbf{Tension par enroulement} & \textbf{Choisir si} \\
\midrule
etoile (Y) & $V_{enr} = V_{composee} / \sqrt{3}$ & $V_{nom\,enr} < V_{composee}$ \\[4pt]
Triangle ($\Delta$) & $V_{enr} = V_{composee}$ & $V_{nom\,enr} = V_{composee}$ \\
\bottomrule
\end{tabular}
\end{center}

\end{tcolorbox}

% -----------------------------------------------------------------------
%  2. FORMULAIRE SYNTHeTIQUE
% -----------------------------------------------------------------------

\begin{tcolorbox}[
  ficheSection=colElec,
  title={2. Formulaire synthetique},
]
\small
\begin{multicols}{2}
\renewcommand{\arraystretch}{1.45}
\begin{tabular}{@{}p{4.5cm}ll@{}}
\toprule
\textbf{Formule} & \textbf{Usage} & \textbf{U.} \\
\midrule
\multicolumn{3}{@{}l}{\footnotesize\textit{Synchronisme}} \\
$N_s = 60f/p$ & Vitesse synchro & tr/min \\
$\omega_m = 2\pi f/p$ & Vitesse angulaire & rad/s \\
$f = p \cdot N/60$ & Frequence imposee & Hz \\
\midrule
\multicolumn{3}{@{}l}{\footnotesize\textit{Circuit equivalent}} \\
$\vec{V}_s = \vec{E} + jX_s\vec{I}_s$ & equation ($R_s=0$) & V \\
$X_s = L_s \omega$ & Reactance sync. & O \\
$E = k\phi\Omega$ & FEM aimants & V \\
$P_{em} = 3V_sE\sin\delta/X_s$ & Puissance em & W \\
$P_{max} = 3V_sE/X_s$ & P max ($\delta=90^\circ$) & W \\
\bottomrule
\end{tabular}

\columnbreak

\renewcommand{\arraystretch}{1.45}
\begin{tabular}{@{}p{4.5cm}ll@{}}
\toprule
\textbf{Formule} & \textbf{Usage} & \textbf{U.} \\
\midrule
\multicolumn{3}{@{}l}{\footnotesize\textit{Couple}} \\
$C_{em} = P_{em}/\Omega$ & Couple em & Nem \\
$C_{nom} = P_{nom}/\Omega_{nom}$ & Couple nominal & Nem \\
\midrule
\multicolumn{3}{@{}l}{\footnotesize\textit{Couplage}} \\
$V_Y = V_{comp}/\sqrt{3}$ & etoile & V \\
$V_\Delta = V_{comp}$ & Triangle & V \\
\midrule
\multicolumn{3}{@{}l}{\footnotesize\textit{Donnees Urbanloop 2024}} \\
$p = 14$ (28 peles) & Urbanloop & - \\
$V_{comp} = 72$ V & Onduleur & V \\
$f_0 = 185{,}6$ Hz & - 60 kmeh$^{-1}$ & Hz \\
\bottomrule
\end{tabular}
\end{multicols}
\normalsize
\end{tcolorbox}

% -----------------------------------------------------------------------
%  3. QUESTIONS TYPES
% -----------------------------------------------------------------------

\begin{tcolorbox}[
  ficheSection=colElec,
  title={3. Questions types - extraites des sujets reels},
]

% -- QT1 -------------------------------------------------------------

% Mini-procedure de resolution type copie concours
\begin{tcolorbox}[questiontype,title={Procedure type copie concours (5 etapes)}]
\begin{enumerate}[leftmargin=1.5em]
  \item \textbf{Qualifier le modele} : hypotheses, domaine de validite, grandeurs d'entree/sortie.
  \item \textbf{Ecrire les lois} : equations de base avec symboles et unites SI.
  \item \textbf{Isoler l'inconnue} : expression litterale, puis application numerique tracee.
  \item \textbf{Controler l'ordre de grandeur} : coherence physique et unite finale.
  \item \textbf{Conclure concours} : phrase courte orientee decision technique.
\end{enumerate}
\end{tcolorbox}

% Visuel de methode : chemin de resolution CAPET
\begin{center}
\begin{tikzpicture}[font=\small, node distance=7mm and 6mm, >=Stealth]
  \node[draw, rounded corners, fill=colElec!8, inner sep=4pt] (n1) {Notions};
  \node[draw, rounded corners, fill=colElec!8, right=of n1, inner sep=4pt] (n2) {Methode};
  \node[draw, rounded corners, fill=colElec!8, right=of n2, inner sep=4pt] (n3) {Application annale};
  \node[draw, rounded corners, fill=colElec!8, right=of n3, inner sep=4pt] (n4) {Verification jury};
  \draw[->, thick, colElec] (n1) -- (n2);
  \draw[->, thick, colElec] (n2) -- (n3);
  \draw[->, thick, colElec] (n3) -- (n4);
\end{tikzpicture}
\end{center}
\begin{tcolorbox}[questiontype,
  title={QT1 - Frequence d'alimentation et couplage - capsule Urbanloop \hfill\normalfont\footnotesize E2 appl. 2024 (Q44-46)}
]
\textbf{enonce.}
La capsule Urbanloop est motorisee par quatre PMSM - 28 peles. La vitesse maximale est $V_0 = 60\,\mathrm{km\!\cdot\!h^{-1}}$, le rayon de roue $R = 0{,}20\,\mathrm{m}$. L'onduleur delivre une tension composee de 72\,V efficace. La tension nominale des enroulements du stator est 72\,V. Calculer la frequence $f_0$ et indiquer le couplage.

\tcbline

\textbf{Corrige.}

\textbf{etape 1 - vitesse de la roue :}
\[
  \omega_{roue} = \frac{V_0}{R} = \frac{60\,000/3600}{0{,}20} = \frac{16{,}67}{0{,}20} = 83{,}3\,\mathrm{rad\!\cdot\!s^{-1}}
\]

\textbf{etape 2 - frequence ($p = 28/2 = 14$) :}
\[
  f_0 = \frac{p\,\omega_{roue}}{2\pi} = \frac{14 \times 83{,}3}{2\pi} \approx \mathbf{185{,}6\,\mathrm{Hz}}
\]

\textbf{etape 3 - couplage :} tension nominale par enroulement = 72\,V = tension composee - \textbf{couplage triangle}. En etoile, chaque enroulement ne recevrait que $72/\sqrt{3} \approx 41{,}6\,\mathrm{V}$.
\end{tcolorbox}

\vspace{4pt}

% -- QT2 -------------------------------------------------------------

\begin{tcolorbox}[questiontype,
  title={QT2 - Servomoteur brushless - definition et difference \hfill\normalfont\footnotesize E2 appl. 2025 (Q18, robots KUKA KR QUANTEC)}
]
\textbf{enonce.}
Les robots KUKA KR QUANTEC sont equipes de servomoteurs brushless. Rappeler la definition d'un servomoteur et preciser la difference avec un moteur classique.

\tcbline

\textbf{Corrige.}

Un \textbf{moteur} convertit une energie electrique en energie mecanique en boucle ouverte : il n'y a pas de retour d'information sur la position ou la vitesse.

Un \textbf{servomoteur} est un moteur associe - un \textbf{capteur de position} (encodeur incremental, resolveur ou capteurs Hall) et - une \textbf{boucle d'asservissement fermee}. Il est commande precisement en position, vitesse ou couple.

Les servomoteurs brushless (PMSM) utilisent des \textbf{aimants permanents au rotor} - pas de collecteur ni de balais - entretien reduit, duree de vie allongee, densite de puissance elevee. La commutation est assuree electroniquement par l'onduleur pilote par les signaux du capteur.
\end{tcolorbox}

\vspace{4pt}

% -- QT3 -------------------------------------------------------------

\begin{tcolorbox}[questiontype,
  title={QT3 - Rendement de la chaene de puissance (robot Bakus) \hfill\normalfont\footnotesize E2 disc. 2023 (Figure 5)}
]
\textbf{enonce.}
La chaene de puissance du robot Bakus comprend : convertisseur statique ($\eta_{CS} = 0{,}97$), machine synchrone ($\eta_{MS} = 0{,}90$) et reducteur ($\eta_{red} = 0{,}72$). La puissance mecanique necessaire est $P_{meca} = 500\,\mathrm{W}$. Exprimer et calculer la puissance electrique $P_{elec}$ fournie par les batteries.

\tcbline

\textbf{Corrige.}

Le rendement global de la chaene est le produit des rendements en serie :
\[
  \eta_{global} = \eta_{CS} \times \eta_{MS} \times \eta_{red} = 0{,}97 \times 0{,}90 \times 0{,}72 = 0{,}630
\]
\[
  P_{elec} = \frac{P_{meca}}{\eta_{global}} = \frac{500}{0{,}630} \approx \mathbf{794\,\mathrm{W}}
\]
Conclusion : les pertes dans la chaene representent $500 \times (1/0{,}63 - 1) \approx 294\,\mathrm{W}$ soit 37 \% de l'energie batt\-erie.
\end{tcolorbox}

\end{tcolorbox}

% -----------------------------------------------------------------------
%  4. ATTENTES DU JURY
% -----------------------------------------------------------------------

\begin{tcolorbox}[
  jury,
  title={4. Attentes du jury - Erreurs recurrentes et points de vigilance},
]

\begin{enumerate}[label=\textbf{\arabic*.}, leftmargin=*, itemsep=5pt]

  \item \textbf{Paires de peles vs nombre de peles.}
  $p$ designe les paires de peles. Une machine - 28 peles a $p = 14$. L'erreur $p = 28$ donne une frequence deux fois trop grande.

  \item \textbf{Couplage etoile/triangle : comparer les tensions.}
  Comparer systematiquement la tension nominale d'un enroulement (donnee constructeur) avec la tension composee disponible. En triangle : $V_{enr} = V_{comp}$. Justifier explicitement.

  \item \textbf{Decrochage de la machine synchrone.}
  - vitesse fixe (reseau), le couple max est atteint pour $\delta = 90^\circ$. Si le couple resistant depasse $C_{max}$, la machine decroche brutalement. En autopilotage (PMSM), ce risque est supprime.

  \item \textbf{Angle de charge et signe.}
  En moteur $\delta > 0$ (rotor en avance sur le champ tournant). En generateur $\delta < 0$. Oublier le signe conduit - une erreur dans le diagramme de Kapp.

  \item \textbf{Servomoteur : mentionner le capteur ET la boucle fermee.}
  Le jury 2025 deplore les reponses qui evoquent le servomoteur sans preciser le capteur de retour (type : resolveur, encodeur) ni le rele de la boucle d'asservissement dans la precision.

  \item \textbf{Rendement de chaene = produit des rendements elementaires.}
  Presenter les trois rendements avant de les multiplier. Ne pas les additionner.

\end{enumerate}

\medskip
\begin{tcolorbox}[
  colback=colJuryFrame!12!colJury,
  colframe=colJuryFrame,
  left=4pt, right=4pt, top=2pt, bottom=2pt,
  fontupper=\small\itshape
]
Rapport jury 2024 : "La partie portant sur l'etude de la machine synchrone (partie 3.3) a ete globalement bien traitee par les candidats qui ont su appliquer la relation entre frequence d'alimentation et vitesse de rotation."\\
Rapport jury 2025 : "Le jury regrette que certains candidats confondent servomoteur et moteur classique, et ne precisent pas la nature de la boucle de retour ni son rele dans la precision du positionnement."
\end{tcolorbox}

\end{tcolorbox}

% -----------------------------------------------------------------------
%  PIED - PReSENCE DANS LES SUJETS
% -----------------------------------------------------------------------

\begin{tcolorbox}[
  enhanced, breakable,
  colback=colGrayLight,
  colframe=colSeparator,
  left=4pt, right=4pt, top=3pt, bottom=3pt,
  before skip=4pt, after skip=0pt,
  fonttitle=\small\bfseries,
  title={Presence dans les sujets analyses (2013-2026)},
]
\begin{center}
\footnotesize\renewcommand{\arraystretch}{1.15}
\begin{tabular}{lccccccccccccc}
\toprule
\textbf{Annee} & 13 & 14 & 15 & 16 & 17 & 18 & 19 & 20 & 21 & 22 & 23 & 24 & 25 \\
\midrule
E1 (commune)   & - & - & - & - & - & - & - & - & - & - & - & - & - \\
E2 discipl. IE & - & - & $\checkmark$ & - & - & - & - & - & - & - & $\checkmark$ & $\checkmark$ & $\checkmark$ \\
\bottomrule
\end{tabular}
\end{center}
\smallskip
\footnotesize\textbf{Sous-themes frequents :}
frequence d'alimentation $f = pN/60$ (3/13) - couplage etoile/triangle (2/13) - servomoteur brushless (2/13) - diagramme de Kapp (2/13) - rendement chaene de puissance (2/13)

\medskip
\textbf{Ordres de grandeur (reperes rapides).}
\begin{itemize}[leftmargin=1.3em]
  \item Reseau BT FR : tension phase-neutre typique \SI{230}{\volt}, tension phase-phase typique \SI{400}{\volt}, frequence \SI{50}{\hertz}.
  \item Courants d'etude frequents en sujet : de quelques \si{\ampere} a quelques dizaines de \si{\ampere}.
  \item En conversion d'energie, un rendement attendu reste generalement entre 0{,}8 et 0{,}98 selon la technologie.
\end{itemize}

\medskip
\textbf{Synthese express (a refaire sans support).}
\begin{enumerate}[leftmargin=1.4em]
  \item Poser les hypotheses du modele et les conventions de signe.
  \item Ecrire 2 a 4 relations cles avec unites.
  \item Realiser un calcul numerique complet et verifier l'ordre de grandeur.
  \item Justifier le resultat par un argument technique de niveau CAPET.
\end{enumerate}

\end{tcolorbox}

\CAPETENDFILE









